% Shared by main.tex and extended/main.tex. Lines marked % CHECK-DATA depend on the data.
\section{Method}
\label{sec:method}

\begin{figure}[t]
  \centering
  \input{\reportdir pipeline.tex}
  \caption{Method. The candidate evaluations (left) each get a decision model (centre): the hidden state $\theta$ with prior $p$ drives the result $x$ through sensitivity $s$ and specificity $t$; the developer sees $x$ and chooses the action $a$; the utility $U$ depends on $\theta$ and $a$ through $B$ and $K$. Two independent prompts elicit the decision parameters (green) and the instrument parameters (purple). Monte Carlo propagates the elicited uncertainty to the value per dollar (right).}
  \label{fig:method}
\end{figure}

\emph{Decision model.} We use a deliberately simple model that gives order-of-magnitude values and rankings, not point predictions (Fig.~\ref{fig:method}). One developer faces one release decision: respond (delay and mitigate) or deploy as planned. A hidden binary state $\theta$ is 1 when the hazardous property is present, so that responding is correct; its prior is $p=P(\theta{=}1)$. One run of the evaluation returns a flag $x\in\{0,1\}$ with sensitivity $s=P(x{=}1\mid\theta{=}1)$ and specificity $t=P(x{=}0\mid\theta{=}0)$. Responding gains $B$ if $\theta=1$ and loses $K$ if $\theta=0$. We value $B$ and $K$ for society, harm avoided and benefits forgone, and assume the developer acts on that value, as published safety frameworks state; the developer's own payoff is an ablation. With stakes $\Lambda=B+K$ and threshold $\pist=K/\Lambda$, acting on a belief $\pi$ is worth $V(\pi)=\Lambda\,[\pi-\pist]^+$ more than deploying, where $[z]^+=\max(z,0)$. A flag arrives with probability $P_1=ps+(1-p)(1-t)$ and moves the belief to $\pi_1=ps/P_1$; a pass moves it to $\pi_0=p(1-s)/(1-P_1)$. The value of one run and its efficiency are
\begin{equation}
\EVSI = P_1\,V(\pi_1)+(1-P_1)\,V(\pi_0)-V(p), \qquad \eta=\EVSI/C, \qquad C=C_\mathrm{build}+C_\mathrm{run},
\label{eq:evsi}
\end{equation}
where $C_\mathrm{build}$ is the cost to build the evaluation from scratch and $C_\mathrm{run}$ the cost of one run against one system. EVSI is zero exactly when no result moves the belief across the threshold: $\pi_0\ge\pist$, so the developer responds regardless, or $\pi_1\le\pist$, so it deploys regardless~\citep{pauker1980threshold}. This happens when the developer is already nearly certain ($p$ far from $\pist$ relative to the evaluation's power), and always when the evaluation is uninformative ($s+t=1$). At fixed stakes, EVSI is largest over thresholds when the developer is exactly undecided ($\pist=p$), where it equals the indifference value $\EVSI^\circ=\Lambda\,p(1-p)\,|s+t-1|$: stakes times the variance of the state times Youden's index $J=s+t-1$~\citep{youden1950index}. A built evaluation informs $n$ decisions over its life (model releases, product versions). Building pays after $n^*=C_\mathrm{build}/(\EVSI-C_\mathrm{run})$ decisions, and never if one run costs more than it is worth.

\emph{Evaluations.} The LLM safety evaluations are those named in LLM system cards. We rank them by how often system cards name them, whether the instrument is public, how checkable the score is and whether one run gives one result, and include those above a fixed score (\inextended{Section~\ref{sec:catalogue}}). They span cyber, biological, loss-of-control, manipulation and societal harms. The physical-AI safety evaluations are published evaluations of embodied systems that score the rate of a physically harmful event. Each has a public score, a cost or validity fact and one build-and-run instrument. They span fidelity levels from text question answering (level 0) to track testing (level 8). Each evaluation gets one release decision written from its sources, a definition of $\theta$, and two paragraphs of sourced facts: one about the decision, one about the evaluation. We invent no number in the model; every number comes from the elicitation.

\emph{Elicitation.} An ensemble of language models stands in for the expert panel that a study of this scale cannot afford. Language models have been used to elicit priors~\citep{capstick2024autoelicit,selby2025hadenough}, following the practice of expert elicitation~\citep{ohagan2006uncertain}, and a diverse LLM crowd forecasts about as well as a human crowd~\citep{schoenegger2024wisdom}. Each evaluation gets two independent prompts. The decision prompt asks for $p$, $B$ and $K$ and never describes the evaluation. The instrument prompt asks for $s$, $t$, $C_\mathrm{build}$, $C_\mathrm{run}$ and $n$ and never states the stakes. For each parameter the member writes its reasoning, then a value it would be surprised to see the truth fall below (5th percentile) or above (95th), then its median. The members are \voiNMembers{} Claude models (\voiMemberList), each answering each prompt twice. % CHECK-DATA: members and repeats of the analysed run
Answers are validated (strict JSON, ordered percentiles, $s>1-t$) and retried once. Of \voiAttempts{} attempts, \voiValidShare{} were valid; \voiInvalidRefusal{} were refusals, all on one biosecurity prompt, and \voiInvalidJson{} were unparseable. % CHECK-DATA The language models are the measurement instrument of this study; this is our main use of generative AI.

\emph{Monte Carlo.} Each elicitation's three percentiles are fitted by least squares with a Beta distribution ($p$, $s$, $t$) or a lognormal ($B$, $K$, $C_\mathrm{build}$, $C_\mathrm{run}$, $n$). The pooled belief per parameter is the equal-weight mixture of the fitted distributions, the linear opinion pool~\citep{stone1961opinionpool,clemen1999combining}. We draw \voiDraws{} parameter vectors per evaluation, with parameters independent, and rank all evaluations on each draw. The central estimate evaluates Eq.~\eqref{eq:evsi} at the median over elicitations of each elicited median.

\emph{Comparing the groups.} The headline metric is $A=P(\eta_\mathrm{phys}>\eta_\mathrm{LLM})$, the probability that a randomly chosen physical-AI evaluation has a higher efficiency than a randomly chosen LLM evaluation, with ties counted as one half. It equals the area under the receiver operating characteristic curve and the Mann-Whitney $U$ statistic divided by the number of pairs~\citep{mcgraw1992cles,vargha2000stochastic}. The percentile curve generalises it to the $q$-th percentile LLM evaluation. We compute both at the central estimate and on every draw. An exact Mann-Whitney test on the central estimates, with p-value $p_\mathrm{MW}$, is the frequentist companion.
